\documentclass[11pt]{article}

\usepackage[margin=1.1in]{geometry}
\usepackage{amsmath,amssymb,amsthm,mathtools}
\usepackage{enumitem}
\usepackage{booktabs}
\usepackage{hyperref}
\usepackage[nameinlink,noabbrev]{cleveref}

\hypersetup{colorlinks=true,linkcolor=blue,citecolor=blue,urlcolor=blue,pdftitle={Topology of the recurrent and transient parameter sets of random walks on abelian groups},pdfauthor={Yitzchak Shmalo}}

\newtheorem{theorem}{Theorem}[section]
\newtheorem{lemma}[theorem]{Lemma}
\newtheorem{proposition}[theorem]{Proposition}
\newtheorem{corollary}[theorem]{Corollary}
\theoremstyle{remark}
\newtheorem{remark}[theorem]{Remark}

\DeclareMathOperator{\supp}{supp}
\DeclareMathOperator{\spn}{span}
\DeclareMathOperator{\conv}{conv}
\DeclareMathOperator{\relint}{relint}
\providecommand{\doi}[1]{\href{https://doi.org/#1}{doi:#1}}
\newcommand{\R}{\mathbb{R}}
\newcommand{\Z}{\mathbb{Z}}
\newcommand{\one}{\mathbf{1}}

\title{Topology of the recurrent and transient parameter sets of random walks on abelian groups}
\author{Yitzchak Shmalo\\
Einstein Institute of Mathematics, The Hebrew University of Jerusalem\\
Givat Ram, Jerusalem, Israel}
\date{}

\begin{document}
\maketitle

\begin{abstract}
Let $R$ be the set of laws, that is, probability vectors, on a finite set $S$ of steps in a finitely generated abelian
group whose homogeneous random walk is recurrent. Li, Mossel and Weiss characterized $R$ and asked when $R$ is simply
connected and when its complement $R^c$ is path connected, simply connected or convex. When no step projects to zero in
the free part of the group, $R$ is homotopy equivalent to a finite bipartite graph whose vertices are certain lines and
planes spanned by projections of steps. This gives the path components of $R$ and their fundamental groups, and $R$ is
simply connected if and only if it is contractible. The complement $R^c$ is empty, contractible, homotopy equivalent to
two points, or homotopy equivalent to a circle, according to the dimension of the span of the projections and whether
the origin is in the relative interior of their convex hull. It is convex exactly when no two or three nonzero
projections are positively dependent.
\end{abstract}

\noindent\textit{Keywords:} random walk, recurrence, abelian group, nerve theorem, homotopy type.\\
\textit{2020 Mathematics Subject Classification:} 60G50, 60J10, 52B05, 55U10.

\section{Introduction}\label{sec:intro}

Let $G$ be a finitely generated abelian group and $S=\{s_1,\dots,s_m\}\subseteq G$ a nonempty finite set. Each point
$a$ of the simplex
\[
A=\Bigl\{a\in\R^m:\ a_i\ge 0,\ \sum_{i=1}^m a_i=1\Bigr\}
\]
defines the homogeneous random walk on $G$ that moves from $g$ to $g+s_i$ with probability $a_i$. We call the points of
$A$ \emph{laws}, and write $\delta_i\in A$ for the point mass at $s_i$; in the examples, $e_1,\dots,e_d$ is the
standard basis of $\Z^d$. Since all states of such a walk have the same type, the walk is either recurrent or transient
even when $S$ does not generate $G$.

Throughout, $G=\Z^n\times H$ with $H$ finite, $\pi\colon G\to\Z^n$ is the projection onto the free part, $v_i=\pi(s_i)$
is the \emph{projection} of the step $s_i$, and
\[
D(a)=\sum_{i=1}^m a_i v_i,\qquad T_a=\{v_i: a_i>0\},
\]
\[
r=\dim\spn\{v_1,\dots,v_m\},\qquad K=\conv\{v_1,\dots,v_m\}.
\]
Distinct steps may have the same projection, and a projection may be zero. We call $D(a)$ the \emph{drift} of the law
$a$, and say that $a$ has \emph{zero drift} if $D(a)=0$. The \emph{support} of a law $a$ is $\supp a=\{i:a_i>0\}$, and $a$ is \emph{supported on} a set
$J$ of indices if $\supp a\subseteq J$. The step set $S$ is \emph{symmetric} if $S=-S$. A set $X$ is \emph{star shaped} about a point
$c\in X$ if it contains the segment from $c$ to each of its points; $c$ is then a \emph{star centre} of $X$.

Li, Mossel and Weiss \cite[Section~6]{LMW} study the set $R\subseteq A$ of recurrent laws, the \emph{recurrent set}, and
its complement $R^c=A\setminus R$, the \emph{transient set}. They prove that $R$ is closed \cite[Theorem~6.2]{LMW}, that $R$ is path
connected when $S$ is symmetric \cite[Theorem~6.3]{LMW}, and they characterize convexity of $R$ \cite[Theorem~6.5,
Corollary~6.6]{LMW}. On the transient side they show that $R^c$ is path connected when some transient law has zero
drift \cite[Theorem~6.9, Corollary~6.10]{LMW}, observe that for $S=\{\pm e_1,\pm e_2\}\subseteq\Z^2$ the set $R^c$ is
path connected but not simply connected, and show that $R^c$ is not convex for symmetric $S$ with a nonzero projection
\cite[Proposition~6.12]{LMW}. After their Example~6.4, a five-step set in $\Z^3$ with $R$ not path connected, they
leave open the classification of path connectedness of $R$ for nonsymmetric $S$, and they pose:
\begin{itemize}[nosep]
\item Question~6.7: when is $R$ simply connected?
\item Question~6.8: when is $R^c$ path connected?
\item Question~6.11: when is $R^c$ simply connected?
\item Question~6.13: when is $R^c$ convex?
\end{itemize}
We answer these questions for every $G$ and every finite $S$. The probabilistic input is the characterization of Li,
Mossel and Weiss.

\begin{proposition}[{\cite[Proposition~6.1]{LMW}}]\label{prop:LMW}
A law $a\in A$ is recurrent if and only if $\dim\spn T_a\le 2$ and $D(a)=0$.
\end{proposition}

Li, Mossel and Weiss obtain the closedness of $R$ from a decomposition into polytopes
\cite[proof of Theorem~6.2]{LMW}. For a linear subspace $W\subseteq\R^n$ let
\[
P_W=\{a\in A:\ D(a)=0,\ a_i=0\ \text{whenever}\ v_i\notin W\}.
\]
By \Cref{prop:LMW}, $R$ is the union of the polytopes $P_W$ over the subspaces $W$ of dimension at most $2$. The
connectedness theorems of \cite{LMW} use more than this decomposition. The proof of \cite[Theorem~6.3]{LMW} joins each
recurrent law by a segment to a law supported on a pair $\{s,-s\}$, which requires $S=-S$, and
\cite[Theorem~6.9]{LMW} assumes a transient law of zero drift, which exists only if $r\ge 3$. Neither argument
detects loops, and loops occur: for $S=\{\pm e_1,\pm e_2,\pm e_3\}\subseteq\Z^3$ the set $R$ is homotopy equivalent
to a circle, and for $S=\{\pm e_1,\pm e_2\}\subseteq\Z^2$ the set $R^c$ is. Both homotopy types can be computed in
general. For $R$, the polytopes satisfy
\[
P_U\cap P_W=P_{U\cap W}\qquad\text{for all subspaces } U,W\subseteq\R^n .
\]
Suppose no $v_i$ is zero. Then $P_{\{0\}}=\emptyset$, and for a line $L$ the polytope $P_L$ is nonempty exactly when
$L$ contains two of the $v_i$ on opposite sides of the origin. Hence, for distinct subspaces $U,W$ of dimension at
most $2$, the intersection $P_U\cap P_W$ is empty or is such a $P_L$. All nonempty intersections are convex, so the
nerve theorem \cite{Bjorner2003} identifies $R$, up to homotopy, with the nerve of this cover, and the nerve
collapses onto the graph whose vertices are these lines and the two-dimensional subspaces of the form $\spn T_a$,
$a\in R$, with an edge for each inclusion of a line in such a subspace (\Cref{thm:rec}). For $R^c$, let $1\le r\le 2$, so that
$R^c=\{a\in A: D(a)\neq 0\}$ by \Cref{prop:LMW}, and suppose that some law with all coordinates positive has zero
drift. Fix a basis of $\spn\{v_1,\dots,v_m\}$, let $\|\cdot\|_\infty$ be the maximum norm in its coordinates and
$\Sigma$ the unit sphere of this norm. There are $\varepsilon>0$ and an embedding $\iota\colon\Sigma\to A$ with
$D(\iota(u))=\varepsilon u$. Put $\rho(a)=\iota\bigl(D(a)/\|D(a)\|_\infty\bigr)$ for $a\in R^c$. Then
\[
D\bigl((1-\lambda)a+\lambda\rho(a)\bigr)=\Bigl(1-\lambda+\frac{\lambda\varepsilon}{\|D(a)\|_\infty}\Bigr)D(a)\neq 0
\qquad(0\le\lambda\le 1),
\]
so the segments from $a$ to $\rho(a)$ lie in $R^c$ and deform $R^c$ onto $\iota(\Sigma)$, which is a pair of points when
$r=1$ and a circle when $r=2$ (\Cref{thm:trans}). In all other cases $R^c$ is empty or star shaped.

We follow \cite{LMW} in
calling the empty set path connected and convex; statements about simple connectedness and contractibility refer to
nonempty spaces, and the empty cases are listed separately.

\subsection{Results on the recurrent set}

A \emph{balanced line} is a one-dimensional subspace $L\subseteq\R^n$ that contains two projections $v_i,v_j$ lying on
opposite sides of the origin. A \emph{plane} is a two-dimensional linear subspace $V\subseteq\R^n$; it is \emph{active}
if some recurrent law $a$ has $\spn T_a=V$. Let $\mathcal B$ be the set of balanced lines, $k=|\mathcal B|$, let
$\mathcal V$ be the set of active planes, write $\mathcal B(V)=\{L\in\mathcal B:L\subseteq V\}$, and let $t$ be the
number of active planes $V$ with $\mathcal B(V)=\emptyset$. All of these sets are finite. The \emph{incidence graph}
$\Gamma$ is the bipartite graph with vertex set $\mathcal B\cup\mathcal V$ and an edge $LV$ whenever $L\subseteq V$.

\begin{theorem}\label{thm:rec}
\begin{enumerate}[label=\textup{(\roman*)},nosep]
\item If $v_i=0$ for some $i$, then $R$ is star shaped about the point mass $\delta_i$, hence contractible.
\item If no $v_i$ is zero, then $R$ is homotopy equivalent to $\Gamma$. Consequently $R$ has exactly
\begin{equation}\label{eq:comp}
t+\min(k,1)
\end{equation}
path components. For each active plane $V$ with $\mathcal B(V)=\emptyset$ the polytope $P_V$ is a convex path component of $R$. If $k>0$, the remaining laws form one path component $R_*$, whose
fundamental group is free of rank
\begin{equation}\label{eq:rank}
\beta=1-k+\sum_{V\in\mathcal V,\ \mathcal B(V)\neq\emptyset}\bigl(|\mathcal B(V)|-1\bigr).
\end{equation}
\end{enumerate}
\end{theorem}

\begin{corollary}[Question~6.7]\label{cor:rec-sc}
Suppose $R\neq\emptyset$. The following are equivalent: $R$ is simply connected; $R$ is contractible; some $v_i$ is zero,
or $t+\min(k,1)=1$ and the balanced lines span a subspace of dimension at most $2$.
\end{corollary}

\begin{corollary}\label{cor:rec-pc}
$R$ is path connected if and only if some $v_i$ is zero or $t+\min(k,1)\le 1$. In particular, if no $v_i$ is zero and
$k>0$, then $R$ is path connected exactly when every active plane contains a balanced line.
\end{corollary}

This follows from \Cref{thm:rec} and \eqref{eq:comp}, and it is the classification asked for in the remark following
\cite[Example~6.4]{LMW}.

For $S=\{\pm e_1,\dots,\pm e_d\}\subseteq\Z^d$ the balanced lines are the $d$ coordinate axes, the active planes are the
$\binom d2$ coordinate planes, and \eqref{eq:rank} gives $\beta=1-d+\binom d2$: the set $R$ is homotopy equivalent to a
wedge of $(d-1)(d-2)/2$ circles, a point for $d\le 2$ and a circle for $d=3$. Adding a step with projection zero makes
$R$ contractible.

\subsection{Results on the transient set}

\begin{theorem}[Questions~6.8 and~6.11]\label{thm:trans}
The homotopy type of $R^c$ is given by the following table, where the relative interior is taken in the affine hull of $K$.
\begin{center}
\begin{tabular}{@{}ll@{}}
\toprule
Case & $R^c$\\
\midrule
$r=0$ & empty\\
$r\ge 3$ & contractible; every law $b$ with $\dim\spn T_b\ge 3$ is a star centre\\
$1\le r\le 2$ and $0\notin\relint K$ & contractible; some point mass $\delta_j$ is a star centre\\
$r=1$ and $0\in\relint K$ & two convex path components\\
$r=2$ and $0\in\relint K$ & contains an embedded circle as a strong deformation retract\\
\bottomrule
\end{tabular}
\end{center}
\end{theorem}

\begin{corollary}\label{cor:trans}
\begin{enumerate}[label=\textup{(\roman*)},nosep]
\item $R^c$ is path connected unless $r=1$ and the nonzero projections do not all lie on the same side of the origin.
\item For $R^c\neq\emptyset$, the following are equivalent: $R^c$ is simply connected; $R^c$ is contractible; $r\ge 3$,
or $1\le r\le 2$ and $0\notin\relint K$. If $r=2$ and $0\in\relint K$, then $\pi_1(R^c)\cong\Z$.
\end{enumerate}
\end{corollary}

Part (i) confirms the belief stated before \cite[Question~6.8]{LMW} that $R^c$ is path connected aside from essentially
one-dimensional examples, and part (ii) bears out their expectation that Question~6.11 has a different
answer.

The condition $0\in\relint K$ holds exactly when some law with all coordinates positive has zero drift
(\Cref{lem:center}). It does not require symmetry: for the steps $(1,0),(0,1),(-2,-3)$ in $\Z^2$ the set $R^c$ has the
homotopy type of a circle. Membership of $0$ in $K$ is not enough: for $(-1,0),(1,0),(0,1)$ the origin lies on an edge
of $K$ and $R^c$ is contractible.

A set of vectors is \emph{positively dependent} if some linear combination of them with all coefficients positive
is zero.

\begin{theorem}[Question~6.13]\label{thm:convex}
Let $Z=\{i:v_i=0\}$ and let $A_Z=\{a\in A:\supp a\subseteq Z\}$, with $A_Z=\emptyset$ when $Z=\emptyset$. The following
are equivalent.
\begin{enumerate}[label=\textup{(\alph*)},nosep]
\item $R^c$ is convex.
\item $R=A_Z$.
\item No two or three of the nonzero projections $v_i$ $(i\notin Z)$, with distinct indices, are positively
dependent.
\end{enumerate}
\end{theorem}

For symmetric $S$ with a nonzero projection, condition (c) fails on the pair $\{v,-v\}$, which recovers
\cite[Proposition~6.12]{LMW}. For the steps $e_1,e_2,e_3,-(e_1+e_2+e_3)$ in $\Z^3$ condition (c) holds, so $R$ is empty
and $R^c=A$ is convex, although the uniform law has zero drift.

\Cref{sec:polytopes} describes the polytopes $P_W$ and their vertices, \Cref{sec:rec} computes the nerve of the cover
of $R$ by them, \Cref{sec:trans} treats $R^c$ by cases on $r$ and on the position of $0$ in $K$, and
\Cref{sec:convex} proves \Cref{thm:convex}.

\section{Zero-drift polytopes}\label{sec:polytopes}

For an index set $J\subseteq\{1,\dots,m\}$ put $Q_J=\{a\in A:\ \supp a\subseteq J,\ D(a)=0\}$, so that
$P_W=Q_J$ for $J=\{i:v_i\in W\}$. Each of these sets is a compact convex polytope, possibly empty.

\begin{lemma}\label{lem:pieces}
\begin{enumerate}[label=\textup{(\roman*)},nosep]
\item $P_U\cap P_W=P_{U\cap W}$ for all subspaces $U,W$.
\item $R=\bigcup P_W$, the union over the finitely many subspaces $W$ of dimension at most $2$ spanned by some of the
$v_i$.
\end{enumerate}
\end{lemma}

\begin{proof}
(i) A law lies in both sets exactly when it has zero drift and is supported on the indices with
$v_i\in U\cap W$. (ii) If $a\in R$, then $a\in P_W$ for $W=\spn T_a$ by \Cref{prop:LMW}. Conversely, if $a\in P_W$ with
$\dim W\le 2$, then $\spn T_a\subseteq W$ and $D(a)=0$, so $a\in R$. Since $P_W$ depends only on the set
$\{i:v_i\in W\}$, $W$ may be replaced by the span of the $v_i$ it contains.
\end{proof}

A \emph{positive circuit} is a nonempty set $C$ of indices, minimal under inclusion among nonempty sets, for which
$\sum_{i\in C}c_iv_i=0$ with all $c_i>0$. A single index with $v_i=0$ is a positive circuit. By minimality the coefficients are unique up to a positive
factor: if $c'$ is a second choice not proportional to $c$, then $c-\lambda c'$, for the largest $\lambda>0$ with
$c-\lambda c'\ge 0$, is nonzero and shows that the projections of a nonempty proper subset of $C$ are positively
dependent. Normalizing the coefficients
gives a law in $Q_C$, the \emph{circuit law} of $C$.

\begin{lemma}\label{lem:vertex}
Let $J$ be an index set. The support of every vertex of $Q_J$ is a positive circuit of size at most
$1+\dim\spn\{v_j:j\in J\}$, and every point of $Q_J$ is a convex combination of the circuit laws of positive circuits
contained in $J$.
\end{lemma}

\begin{proof}
Let $a$ be a vertex of $Q_J$ with support $C$. If the vectors $(1,v_i)\in\R^{1+n}$, $i\in C$, were linearly dependent, a
nonzero $y$ supported on $C$ with $\sum y_i=0$ and $\sum y_iv_i=0$ would give $a\pm\theta y\in Q_J$ for small
$\theta>0$, contradicting that $a$ is a vertex. Hence these vectors are independent; they lie in
$\R\times\spn\{v_j:j\in J\}$, so $|C|\le 1+\dim\spn\{v_j:j\in J\}$. If the projections of a nonempty proper subset
$C'\subsetneq C$ were positively dependent, the normalized law of a vanishing positive combination $a'$ would satisfy
$\sum_{i\in C}(a_i-a'_i)(1,v_i)=0$ with $a-a'\neq 0$, again a dependence. So $C$ is a positive circuit and $a$ is its circuit
law. The last assertion holds because a compact convex polytope is the convex hull of its vertices.
\end{proof}

\begin{lemma}\label{lem:lines}
Suppose no $v_i$ is zero.
\begin{enumerate}[label=\textup{(\roman*)},nosep]
\item $P_{\{0\}}=\emptyset$, and for a line $L$ the polytope $P_L$ is nonempty if and only if $L$ is balanced.
\item A plane $V$ is active if and only if it contains the projections $v_i$, $i\in C$, of a positive circuit $C$ of
size $3$ whose projections span $V$, or two distinct balanced lines.
\item If a plane $V$ is not active, then $P_V=\bigcup_{L\in\mathcal B(V)}P_L$.
\item If $L,L'$ are distinct balanced lines, the plane $L+L'$ is active.
\end{enumerate}
\end{lemma}

\begin{proof}
(i) No law has $T_a\subseteq\{0\}$, so $P_{\{0\}}=\emptyset$. Write the projections on $L$ as $c_i\omega$ with $\omega\neq 0$; a zero-drift law on them exists if
and only if both signs of $c_i$ occur.

(ii) Let $a\in P_V$ with $\spn T_a=V$. By \Cref{lem:vertex} applied to $J=\supp a$, $a$ is a convex combination with
positive weights of circuit laws of positive circuits in $\supp a$, of size at most $3$, whose union is $\supp a$. A
positive circuit of size $3$ cannot lie on a line, since among three positively dependent collinear nonzero vectors two are
negative multiples of each other and already form a positive circuit. Positive circuits of size $2$ are pairs of
projections that are negative multiples of each other, and they lie on balanced lines. If no circuit has size $3$ and all circuits lie on one line, then
$\spn T_a$ is that line, a contradiction. Conversely, a circuit law of three projections spanning $V$ lies in $P_V$ and
has $\spn T_a=V$, and so does the midpoint of laws in $P_L$ and $P_{L'}$ for distinct balanced lines $L,L'\subseteq V$.

(iii) If $V$ is not active and $a\in P_V$, then $\spn T_a$ is a line $L\subseteq V$ (it is not $\{0\}$, by (i)), and
$a\in P_L$ forces $L$ to be balanced. (iv) follows from the last sentence of the proof of (ii).
\end{proof}

\section{The recurrent set}\label{sec:rec}

\begin{proof}[Proof of \Cref{thm:rec}]
(i) Let $v_i=0$ and $a\in R$. For $0\le\lambda\le 1$ the law $(1-\lambda)a+\lambda \delta_i$ has drift $(1-\lambda)D(a)=0$
and $\spn T_{(1-\lambda)a+\lambda \delta_i}\subseteq\spn T_a$, since adding the zero vector does not enlarge a span. By
\Cref{prop:LMW} the whole segment lies in $R$.

(ii) Assume no $v_i$ is zero. By \Cref{lem:pieces}(ii) and \Cref{lem:lines}(i),(iii), the sets $P_L$ $(L\in\mathcal B)$
and $P_V$ $(V\in\mathcal V)$ are nonempty and cover $R$. We determine which subfamilies have a common point, using
\Cref{lem:pieces}(i) and \Cref{lem:lines}(i). For distinct balanced lines $P_L\cap P_M=P_{\{0\}}=\emptyset$. For a
balanced line and an active plane, $P_L\cap P_V=P_L$ if $L\subseteq V$ and $P_L\cap P_V=P_{\{0\}}=\emptyset$ otherwise.
Two or more distinct planes meet in a subspace of dimension at most one, so a family of two or more active planes has a
common point exactly when all of them contain one balanced line $L$, and the common intersection is then $P_L$. Hence
the nerve $N$ of the cover is the union of the full simplices
\[
\Delta_L\quad\text{on the vertex set}\quad\{L\}\cup\{V\in\mathcal V:L\subseteq V\},\qquad L\in\mathcal B,
\]
together with one isolated vertex for every active plane $V$ with $\mathcal B(V)=\emptyset$. If $L\neq L'$, then
$\Delta_L\cap\Delta_{L'}$ is the single vertex $L+L'$, which is active by \Cref{lem:lines}(iv): a common face can contain
only planes containing both lines.

All members of the cover and all their nonempty intersections are compact convex sets, hence contractible. The sign
conditions on the coordinates $a_i$ and on the coordinates of $D(a)$ cut $A$ into a polytopal complex, a regular
CW complex in which every $P_W$ is a subcomplex. Call two members of the cover linked if they are joined by a chain
of members in which consecutive ones meet. The unions of the classes of linked members are connected, pairwise
disjoint and closed, so they are the path components of $R$, and they correspond to the components of $N$. The
nerve theorem of Borsuk, in the form of \cite[Theorem~6 and Remark~7]{Bjorner2003}
for a connected regular CW complex covered by subcomplexes whose nonempty finite intersections are contractible,
applies to each of them. Hence $R$ is homotopy equivalent to $N$; we write $\simeq$ for homotopy equivalence.

Write $\Delta_L=L*F_L$ for the cone with apex $L$ over $F_L$, the simplex on the planes containing $L$. A face $\sigma$ of $F_L$ of dimension
at least one contains two planes, whose intersection is $L$, so neither $\sigma$ nor $L*\sigma$ lies in any other
$\Delta_{L'}$. Remove the pairs $(\sigma,L*\sigma)$, for the faces $\sigma$ of $F_L$ of dimension at least one, in order of
decreasing dimension. When $\sigma$ is removed, every face strictly containing it other than $L*\sigma$ has already been
removed, so $\sigma$ is a free face of $L*\sigma$ and the removal is an elementary collapse. Doing this for every $L$
leaves exactly the graph $\Gamma$. An elementary collapse is a deformation retraction, because a simplex deformation
retracts onto the union of all its facets but one, so $R\simeq N\simeq\Gamma$.

The graph $\Gamma$ has one isolated vertex for each active plane with $\mathcal B(V)=\emptyset$. If $k>0$, every other
vertex lies in one path component: two line vertices $L,L'$ are joined through $L+L'$, and a plane vertex with
$\mathcal B(V)\neq\emptyset$ is adjacent to a line. This gives \eqref{eq:comp}, since a homotopy equivalence induces a
bijection of path components. The path component of $R$ that corresponds to an isolated vertex $V$ is $P_V$, since
$P_V$ meets no other member of the cover, and $P_V$ is convex. The path component of $\Gamma$ containing the lines has
$k+|\mathcal V|-t$ vertices and $\sum_V|\mathcal B(V)|$ edges. A connected graph with $n_0$ vertices and $n_1$ edges is homotopy equivalent to a
wedge of $n_1-n_0+1$ circles and has free fundamental group of that rank \cite[Proposition~1A.2 and its proof]{Hatcher}, which gives \eqref{eq:rank}.
\end{proof}

\begin{proof}[Proof of \Cref{cor:rec-sc}]
A contractible space is simply connected. If some $v_i$ is zero, $R$ is contractible by \Cref{thm:rec}(i). Assume no
$v_i$ is zero. A simply connected space is path connected, so \eqref{eq:comp} equals one. Then $R\simeq\Gamma$ with
$\Gamma$ connected, or $R$ is a
single convex polytope $P_V$. It remains to show that, when $R$ is path connected and $k>0$, the graph $\Gamma$ is a tree
if and only if the balanced lines span at most two dimensions; here a tree is a finite connected graph containing no embedded circle, and for a
finite connected graph, being simply connected, being contractible and being a tree are equivalent
\cite[Section~1.A, Proposition~1A.2]{Hatcher}. If the balanced lines span one dimension, then $k=1$
and $\Gamma$ is a star. If they span a plane $W_0$, then $k\ge 2$, $W_0$ is active by \Cref{lem:lines}(iv), and $W_0$ is
adjacent to every line vertex; any other active plane contains at most one balanced line, since two would span $W_0$, so it
is a leaf, and $\Gamma$ is a tree. If the balanced lines span at least three dimensions, choose linearly independent
$L_1,L_2,L_3\in\mathcal B$. The planes $L_1+L_2$, $L_2+L_3$, $L_1+L_3$ are distinct and active, and
$L_1,\ L_1+L_2,\ L_2,\ L_2+L_3,\ L_3,\ L_1+L_3$ is a closed path in $\Gamma$ without repeated vertices, so
$\Gamma$ contains an embedded circle and is not a tree.
\end{proof}

If $S$ is symmetric and no $v_i$ is zero, let $V$ be an active plane. By \Cref{lem:lines}(ii), $V$ contains two
balanced lines or the projections of a positive circuit $C$ of size $3$. In the second case pick $i\in C$; since
$-s_i\in S$ and $\pi(-s_i)=-v_i\neq 0$, the line through $v_i$ is balanced and lies in $V$. Hence $t=0$, and $R$ is
path connected by \Cref{cor:rec-pc}. If some $v_i$ is zero, $R$ is contractible by \Cref{thm:rec}(i). This recovers
\cite[Theorem~6.3]{LMW}.

\begin{remark}\label{rem:minimal}
Li, Mossel and Weiss state that their Example~6.4 is a minimal example of a step set whose $R$ is not path connected, with $n=3$ and $|S|=5$.
This follows from \Cref{cor:rec-pc}: if $R$ is not path connected, there are at least five steps and $r\ge 3$. Indeed, by
\Cref{cor:rec-pc} no $v_i$ is zero and either $t\ge 2$, or $t\ge 1$ and $k\ge 1$. An active plane $V$ with
$\mathcal B(V)=\emptyset$ contains the projections of a positive circuit $C_V$ of size $3$, by
\Cref{lem:lines}(ii). If $L$ is a balanced line not contained in $V$, the two indices whose projections lie on $L$ on
opposite sides of the origin are not in $C_V$. This gives five indices, hence five steps, whose projections span three
dimensions. If $U\neq V$ are two such planes, with circuits $C_U$ and $C_V$, then at most one index of each circuit has
its projection on the subspace $U\cap V$, which has dimension at most one. Indeed, two such projections would either
be negative multiples of each other, and then form a smaller positive circuit, or positive multiples of each other,
and then the positive dependence would put the third projection on the same line. So $C_U$ and $C_V$ share at most one index, and again there are at least
five steps, whose projections span $U+V$.
\end{remark}

\section{The transient set}\label{sec:trans}

If $r=0$, every $v_i$ is zero and $R=A$ by \Cref{prop:LMW}, so $R^c$ is empty. From now on in this section
$r\ge 1$. Let $V_0=\spn\{v_i\}$, fix a linear isomorphism $V_0\cong\R^r$, let $w_i\in\R^r$ be
the image of $v_i$, and let $\Phi$ be the $r\times m$ matrix with columns $w_i$, so that $D(a)=0$ if and only if $\Phi a=0$.
When $0\in K$, the affine hull of $K$ is the linear space $V_0$, so when $0\in K$ the relative interior of $K$ is its
interior in $V_0$.

\begin{lemma}\label{lem:sep}
If $0\notin\relint K$, there is a linear functional $\ell$ on $V_0$ with $\ell(v_i)\ge 0$ for all $i$ and
$\ell(v_j)>0$ for some $j$.
\end{lemma}

\begin{proof}
Let $\langle\cdot,\cdot\rangle$ be the standard inner product of $\R^n$ and $|\cdot|$ its norm. If $0\notin K$, let $p$
be the point of $K$ nearest to the origin; for $x\in K$ the function $\lambda\mapsto|p+\lambda(x-p)|^2$ is minimal at $\lambda=0$ on $[0,1]$, so $\langle p,x-p\rangle\ge 0$ and
$\langle p,x\rangle\ge|p|^2>0$. Take $\ell=\langle p,\cdot\rangle$. Otherwise $0$ lies on the relative boundary of $K$
in $V_0$. Choose $z_\nu\in V_0\setminus K$ with $z_\nu\to 0$, let $p_\nu$ be the point of $K$ nearest to $z_\nu$, and
$\xi_\nu=(p_\nu-z_\nu)/|p_\nu-z_\nu|$. As before $\langle \xi_\nu,x-p_\nu\rangle\ge 0$ for $x\in K$, and
$|p_\nu|\le 2|z_\nu|\to 0$. A subsequence of $\xi_\nu$ converges to a unit vector $\xi\in V_0$ with
$\langle \xi,x\rangle\ge 0$ on $K$; take $\ell=\langle \xi,\cdot\rangle$. In both cases $\ell\neq 0$ on $V_0$, and since the
$v_i$ span $V_0$ some $\ell(v_j)$ is nonzero, hence positive.
\end{proof}

\begin{lemma}\label{lem:center}
The origin lies in $\relint K$ if and only if there is $q\in A$ with all $q_i>0$ and $\Phi q=0$. In that case there is a
real $m\times r$ matrix $B$ with $\one^{\mathsf T}B=0$ and $\Phi B=I_r$, where $\one=(1,\dots,1)^{\mathsf T}\in\R^m$.
\end{lemma}

\begin{proof}
Suppose $0\in\relint K$. Then $0\in K$, so the affine hull of $K$ is $V_0$. For each $i$ with $v_i\neq 0$ choose
$\eta_i>0$ with $-\eta_iv_i\in K$, write $-\eta_iv_i=D(b^{(i)})$ with $b^{(i)}\in A$, and put
$q^{(i)}=(\eta_i\delta_i+b^{(i)})/(1+\eta_i)$; for $v_i=0$ put $q^{(i)}=\delta_i$. Each $q^{(i)}$ has zero drift and positive
$i$th coordinate, so their average $q$ has all coordinates positive and $\Phi q=0$.

Conversely, let $q$ be as stated. Given $x\in\R^r$, choose $y\in\R^m$ with $\Phi y=x$, which is possible because the $w_i$
span $\R^r$. Then $y-(\one^{\mathsf T}y)q$ still maps to $x$ and has coordinate sum zero. Taking these preimages for the
standard basis vectors as the columns of $B$ gives $\one^{\mathsf T}B=0$ and $\Phi B=I_r$. For $x\in\R^r$ the vector
$q+Bx$ has coordinate sum one and is positive for $x$ small, and $\Phi (q+Bx)=x$. Hence a neighbourhood of $0$ in $V_0$ lies
in $K$, that is, $0\in\relint K$.
\end{proof}

\begin{proof}[Proof of \Cref{thm:trans}]
The case $r=0$ was settled at the start of this section.

Let $r\ge 3$ and let $b\in A$ satisfy $\dim\spn T_b\ge 3$; the uniform law is one choice. For $a\in R^c$ consider the
segment $h(a,\lambda)=(1-\lambda)a+\lambda b$, $0\le\lambda\le 1$. For $\lambda>0$ the support of $h(a,\lambda)$ contains
the support of $b$, so $\dim\spn T_{h(a,\lambda)}\ge 3$ and $h(a,\lambda)\in R^c$ by \Cref{prop:LMW}, whatever the
drift. Thus $R^c$ is star shaped about $b$.

Let $1\le r\le 2$. Then $\dim\spn T_a\le 2$ for every $a$, and \Cref{prop:LMW} gives
\begin{equation}\label{eq:Tc}
R^c=\{a\in A:\ \Phi a\neq 0\}.
\end{equation}
If $0\notin\relint K$, take $\ell$ and $j$ from \Cref{lem:sep}. For $a\in R^c$ and $0<\lambda\le 1$,
\[
\ell\bigl(D((1-\lambda)a+\lambda \delta_j)\bigr)=(1-\lambda)\ell(D(a))+\lambda\ell(v_j)>0,
\]
since $\ell(D(a))\ge 0$. By \eqref{eq:Tc} the segment from $a$ to $\delta_j$ lies in $R^c$, so $R^c$ is star shaped about
$\delta_j$.

Finally let $0\in\relint K$, and take $q$ and $B$ from \Cref{lem:center}. Put
\[
q_*=\min_iq_i>0,\qquad \kappa=\max_i\sum_{j=1}^r|B_{ij}|,\qquad \varepsilon=\frac{q_*}{2\kappa},
\]
where $\kappa>0$ because $\Phi B=I_r$. Let $\Sigma=\{u\in\R^r:\|u\|_\infty=1\}$, a sphere of dimension $r-1$, and
$\iota(u)=q+\varepsilon Bu$. For $\|u\|_\infty\le 1$ the coordinates of $\iota(u)$ sum to one and satisfy
\begin{equation}\label{eq:sphere}
\iota(u)_i\ge q_i-\varepsilon\kappa\ge q_*/2>0,\qquad \Phi \iota(u)=\varepsilon u .
\end{equation}
Hence $\iota(\Sigma)\subseteq R^c$, and $\iota$ maps $\Sigma$ homeomorphically onto its image, with inverse $a\mapsto \Phi a/\varepsilon$.
For $a\in R^c$ define
\[
u(a)=\frac{\Phi a}{\|\Phi a\|_\infty},\qquad \rho(a)=\iota(u(a)),\qquad \tilde h(a,\lambda)=(1-\lambda)a+\lambda\rho(a),
\]
which are continuous on $R^c$ by \eqref{eq:Tc}. By convexity of $A$, $\tilde h(a,\lambda)\in A$, and
\[
\Phi \tilde h(a,\lambda)=\bigl((1-\lambda)\|\Phi a\|_\infty+\lambda\varepsilon\bigr)u(a)\neq 0,
\]
so $\tilde h$ stays in $R^c$. It is the identity at $\lambda=0$ and maps onto $\iota(\Sigma)$ at $\lambda=1$. If $a=\iota(u)$ with
$u\in\Sigma$, then $\Phi a=\varepsilon u$, so $u(a)=u$ and $\rho(a)=a$, and $\tilde h(a,\lambda)=a$ for all $\lambda$. Thus $\iota(\Sigma)$
is a strong deformation retract of $R^c$; it is two points for $r=1$ and a circle for $r=2$. For $r=1$ write
$w_i=c_i\in\R$: both signs occur, and by \eqref{eq:Tc} the set $R^c$ is the disjoint union of the nonempty convex open sets
$\{\sum a_ic_i>0\}$ and $\{\sum a_ic_i<0\}$.
\end{proof}

\begin{proof}[Proof of \Cref{cor:trans}]
By \Cref{thm:trans}, $R^c$ fails to be path connected only when $r=1$ and $0\in\relint K$. When $r=1$, $K$ is a segment
on the line $V_0$, and $0$ lies in its relative interior exactly when nonzero projections occur on both sides of the
origin. Part (ii) follows from the table, since two points are not path connected and $\pi_1$ of a circle is $\Z$
\cite[Theorem~1.7]{Hatcher}.
\end{proof}

For $S=\{\pm e_1,\pm e_2\}$ the origin is the centre of the square $K$, so $R^c$ deformation retracts onto a circle, in
agreement with the observation of \cite{LMW} that this $R^c$ is path connected and not simply connected. For $r\ge 3$, $R^c$ is contractible even when it contains zero-drift laws, which is the
situation of \cite[Theorem~6.9]{LMW}, and also when it does not: for $S=\{e_1,e_2,e_3\}$ every law has nonzero drift.

\section{Convexity of the transient set}\label{sec:convex}

\begin{proof}[Proof of \Cref{thm:convex}]
(c)$\Rightarrow$(b). A law in $A_Z$ has zero drift and $T_a\subseteq\{0\}$, so $A_Z\subseteq R$ by
\Cref{prop:LMW}. Let $a\in R\setminus A_Z$, and let $a'$
be $a$ conditioned on the complement of $Z$. Since the steps with projection zero do not contribute to the drift, $D(a')=0$, and
$\spn T_{a'}\subseteq\spn T_a$ has dimension at most two. By \Cref{lem:vertex} the polytope $Q_{\supp a'}$ has a vertex whose
support is a positive circuit of size at most $3$; it is not a single index because the projections on the
support of $a'$ are nonzero. This contradicts (c), so $R=A_Z$.

(b)$\Rightarrow$(a). If $R=A_Z$, then $R^c=\{a\in A:\sum_{i\notin Z}a_i>0\}$, which is convex (and empty if
$Z=\{1,\dots,m\}$).

(a)$\Rightarrow$(c). Suppose $C$ is a set of two or three indices outside $Z$ with $\sum_{i\in C}c_iv_i=0$, $c_i>0$. We may
take $C$ minimal, so it is a positive circuit; let $\mu$ be its circuit law. Its support spans at most two dimensions,
so $\mu\in R$. Fix $i\in C$ and let $b$ be $\mu$ conditioned on $C\setminus\{i\}$. The point mass $\delta_i$ has drift $v_i\neq 0$,
and $b$ has drift $-\mu_iv_i/(1-\mu_i)\neq 0$, so both are transient, while $\mu=\mu_i\delta_i+(1-\mu_i)b$ is recurrent. Hence $R^c$
is not convex.
\end{proof}

\section*{Acknowledgments and funding}

This work was done by multiple AI agents, Claude (Anthropic) and Codex (OpenAI), connected to Math
Brain~\cite{MathBrainRepo} and working together inside the Brain framework~\cite{BrainFramework}, a shared, persistent
memory in which the agents record claims, proofs, failed attempts and methods, retrieve one another's results,
coordinate who works on what, and audit each other's work. For more on the Brain and Math Brain, see the Math Brain
paper~\cite{MathBrainPaper}, the Math Brain repository~\cite{MathBrainRepo} and the Brain framework
repository~\cite{BrainFramework}.

The research presented in this paper was supported by the European Research Council (ERC) under the European
Union's Horizon Europe research and innovation programme (grant agreement No.\ 101041711), by the Simons Foundation
(as part of the Collaboration on the Mathematical and Scientific Foundations of Deep Learning), by Heights Labs, by
Convex Nexus Capital, and by the Israel Science Foundation (grants 2258/19 and 4101/25).

\section*{Declaration of competing interest}

The author holds equity in Heights Labs and Convex Nexus Capital and has a management role at Convex Nexus Capital.

\section*{Declaration of generative AI and AI-assisted technologies}

This paper is part of a broader research project on proving open problems with AI systems.
The arguments of \Cref{sec:polytopes,sec:rec,sec:trans,sec:convex} were developed with AI systems (OpenAI Codex and
Anthropic Claude agents), which also drafted the text and wrote the script
\texttt{check\_lmw.py}. They build on the work of Li, Mossel and Weiss \cite{LMW}: \Cref{prop:LMW} is the probabilistic
input, the decomposition of $R$ into the polytopes $P_W$ comes from the proof of \cite[Theorem~6.2]{LMW}, and the
questions answered here are theirs. The topological step is the nerve theorem of Borsuk, in the form
given by Bj\"orner \cite{Bjorner2003}. The author directed
the work and is responsible for the content of the paper and for decisions concerning its dissemination.

\section*{Data and code availability}

The script \texttt{check\_lmw.py}, which evaluates \eqref{eq:comp}, \eqref{eq:rank} and the criteria of
\Cref{thm:trans,thm:convex} on explicit step sets in exact rational arithmetic, is in the repository
\begin{center}\url{https://github.com/yspennstate/recurrent-transient-parameter-sets}.\end{center}

\appendix
\section{Notation}

\begin{center}
\small
\begin{tabular}{@{}lp{0.52\textwidth}l@{}}
Symbol & Meaning & First use\\[2pt]
$G=\Z^n\times H$ & finitely generated abelian group, $H$ finite & \Cref{sec:intro}\\
$S=\{s_1,\dots,s_m\}$ & finite step set & \Cref{sec:intro}\\
$A$, $\delta_i$ & simplex of laws on $S$, point mass at $s_i$ & \Cref{sec:intro}\\
$\supp a$ & support $\{i:a_i>0\}$ of a law & \Cref{sec:intro}\\
$R$, $R^c$ & recurrent set, transient set & \Cref{sec:intro}\\
$\pi$, $v_i$ & projection to $\Z^n$, projections of the steps & \Cref{sec:intro}\\
$D(a)$, $T_a$ & drift of $a$, set of the $v_i$ with $a_i>0$ & \Cref{sec:intro}\\
$r$, $K$ & dimension of $\spn\{v_i\}$, convex hull of the $v_i$ & \Cref{sec:intro}\\
$X$ & a star-shaped set & \Cref{sec:intro}\\
$\|\cdot\|_\infty$, $\Sigma$ & maximum norm, its unit sphere & \Cref{sec:intro}\\
$\varepsilon$, $\iota$, $\rho$, $\lambda$ & radius and embedding of $\Sigma$ in $R^c$, retraction of $R^c$, homotopy parameter & \Cref{sec:intro}\\
$L$, $M$, $\mathcal B$, $k$ & lines, balanced lines and their number & \Cref{sec:intro}, \Cref{sec:rec}\\
$V$, $\mathcal V$, $\mathcal B(V)$, $t$ & a plane, active planes, balanced lines in $V$, number with $\mathcal B(V)=\emptyset$ & \Cref{sec:intro}\\
$\Gamma$ & incidence graph of lines and planes & \Cref{sec:intro}\\
$R_*$, $\beta$ & path component of $R$ containing the $P_L$, rank of $\pi_1(R_*)$ & \Cref{thm:rec}\\
$\simeq$ & homotopy equivalence & \Cref{sec:rec}\\
$Z$, $A_Z$ & indices of steps with projection zero, laws supported on them & \Cref{thm:convex}\\
$U$, $W$, $P_W$, $Q_J$ & linear subspaces, zero-drift laws supported on the indices with $v_i\in W$, on $J$ & \Cref{sec:intro}, \Cref{sec:polytopes}\\
$C$, $c_i$ & positive circuit, its coefficients & \Cref{sec:polytopes}\\
$N$, $\Delta_L$, $F_L$, $\sigma$ & nerve of the cover, its full simplices, the simplex on the planes containing $L$, a face of $F_L$ & \Cref{sec:rec}\\
$V_0$, $w_i$, $\Phi$ & span of the $v_i$, coordinates in $\R^r$, matrix of the $w_i$ & \Cref{sec:trans}\\
$\ell$ & linear functional separating $K$ from the origin & \Cref{lem:sep}\\
$\langle\cdot,\cdot\rangle$, $|\cdot|$, $\xi$ & standard inner product and norm of $\R^n$, unit normal & proof of \Cref{lem:sep}\\
$p$, $z_\nu$, $p_\nu$ & nearest point of $K$ to the origin; points outside $K$ tending to $0$, the points of $K$ nearest to them & proof of \Cref{lem:sep}\\
$\one$, ${}^{\mathsf T}$, $I_r$ & all-ones vector in $\R^m$, transpose, identity matrix & \Cref{lem:center}\\
$b$, $b^{(i)}$, $\eta_i$ & a law of $A$; laws with $D(b^{(i)})=-\eta_iv_i$ and the positive scalars $\eta_i$ & \Cref{thm:trans}, \Cref{lem:center}\\
$h(a,\lambda)$, $\tilde h(a,\lambda)$ & the segments $(1-\lambda)a+\lambda b$ and $(1-\lambda)a+\lambda\rho(a)$ & proof of \Cref{thm:trans}\\
$q$, $q_*$, $B$, $\kappa$, $u(a)$ & data of the deformation retraction & \Cref{lem:center}, \Cref{thm:trans}\\
$\mu$ & circuit law of a positive circuit & proof of \Cref{thm:convex}\\
\end{tabular}
\end{center}

\noindent Terms, with the place of their definition: law, drift, support, supported on, symmetric, star shaped, star
centre (\Cref{sec:intro}); recurrent set, transient set (\Cref{sec:intro}); balanced line, plane, active plane,
incidence graph (\Cref{sec:intro}, before \Cref{thm:rec}); positively dependent (before \Cref{thm:convex});
positive circuit, circuit law (\Cref{sec:polytopes}).

\begin{thebibliography}{99}

\bibitem{Bjorner2003}
A.~Bj\"orner, \emph{Nerves, fibers and homotopy groups}, J.\ Combin.\ Theory Ser.~A \textbf{102} (2003), 88--93.
\doi{10.1016/S0097-3165(03)00015-3}.

\bibitem{Hatcher}
A.~Hatcher, \emph{Algebraic Topology}, Cambridge University Press, Cambridge, 2002.

\bibitem{LMW}
R.~Li, E.~Mossel and B.~Weiss, \emph{Monotonicity, topology, and convexity of recurrence in random walks}, Ann.\ Appl.\
Probab.\ \textbf{36} (2026), 4291--4310. \doi{10.1214/26-AAP2329}; arXiv:2403.04108.

\bibitem{MathBrainRepo}
Y.~Shmalo, \emph{Math Brain: public snapshot, code and tests}, 2026.
\url{https://github.com/yspennstate/math-brain-public}

\bibitem{MathBrainPaper}
Y.~Shmalo, \emph{Math Brain: shared, auditable memory for multi-agent mathematical research}, October 2026.
\url{https://github.com/yspennstate/math-brain-public/tree/main/papers/math-brain-2026-10-24}

\bibitem{BrainFramework}
Y.~Shmalo, \emph{Brain Framework}, public software repository, 2026.
\url{https://github.com/yspennstate/brain-framework-public}

\end{thebibliography}

\end{document}
